
\comment{Conventions/notations:
\begin{itemize}
\item not using todo-notes, instead using bespoke macros to easily distinguish different types of todos, and so they are not visually overbearing.
    \item Now using a macro for arrow-type (currently $\arrow$). \anupam{I sometimes forget and use $\to$, not a big deal anymore.} No dot after binders in grammars (it is a meta-level notation). $\lambda$ etc.\ binds as strongly as possible (unless there is a dot).
    \item application always without brackets, i.e.\ $st$ not $s(t)$. This is consistent with above convention.
    \item $\closure A$ is the closure of a set $A$ of (co)terms under untyped application. Just saying `typed' for typed version at the moment.
    \item Notation convention: everything to do with $\muLJ$ and the iteration rules have standard notations (e.g.\ $\muLJ$, $\lr \mu$, $\iter{}$); everything to do with $\cmuLJ$ and left-unfolding rules currently have `primed' notations (e.g.\ $\muLJnorec$, $\mulunf$, $\iter{}'$). \anupam{I am not necessarily sold on this notation, I welcome alternative suggestions. If we keep might be better to call $\cmuLJ$ something like $\muLJnorec_{\mathrm{reg}}$}
    \item do not say `number-theoretic function', better to say `function on natural numbers' or even `numerical function'.
    \item use $\Sigma, \Gamma$ for lists/cedents of types to avoid too many arrows.
    \item use $\seqar$ for sequent arrow, not $\vdash$ since that is reserved for provability (mathematical judgement, not syntactic delimiter).
    \item Using mathrm for meta-level sets (e.g. $\var, \Var, \Types$).
    \item using mathsf for term constants, proof rules etc.
    \item Do not use cref: does not work with this document class
\end{itemize}
}